\head{第五卷增补合计核查}
\label{rep:suppIII}
本章接续原第十七版报告的最后范围，以实际合订的第50--68批归档为准。底本为合并英译的物理第66页交换图至第77页末；法文扫描的扫描顺序及印刷页码并不总连续，尤其15.2之后依次为物理160、159、158页（原稿116、117、118）。以下按现行合订译文报告，不把作者工作笔记提升为正式EGA刊行正文。正式EGA第四章章首不在本批新增范围；其断点仍是 \texttt{EGAIV.1.pdf} 物理221页提要、222--223页章首与223页（1.1.1）。
\section*{第一辑：英译物理第66--68页}
\paragraph{第50--51批：两线丛、交模和交换图。}
按英文续页补全PS末句之后的交换图，注意$a,b$的整体符号随行列式自然同构的选择而变，不能擅填。由$a,b$生成的是理想，$V(a),V(b)$是零概形；若记理想为$I_a,I_b$，概形交对应$I_a+I_b$，概形并对应$I_a\cap I_b$，均不是与原交模直接同类型的对象。主部记号$P_v$的阶数从前文读作一阶。只按同一选择解释底本所说的交换性，不把残存符号的歧义当已解决。
\paragraph{第52批：非超奇异层的序列。}
英译与法文物理102页均存第二条正合列的靶秩问题。正文保留原显示式；译注指出交模必须用前文带对偶的一阶主部模，而且由前图推得的商需要去掉共同线。这个按秩与交模重建的候选靶不是原稿写成的改式。扭曲线丛使用张量乘，不能把英译乘号直接解释为直积。
\paragraph{第53--54批：可逆模换元、普通奇异零点。}
法文物理102页动词是“应当总结”，不是已有命题。只有一阶项已消失时，更换局部生成元的高阶项才按所述单位倍数变换；所谓孤立性的泛化须排除原点。法、英关于正规性维数的词句均存疑；若用 Cohen--Macaulay 与穿孔谱正则推出正规，维数门槛应施于超曲面局部环，而不是不经说明施于环境。
\paragraph{第55批：类比图及划去稿的分层。}
重画线丛截面的图，不把斜箭头另当成单射或满射；指数$K^{-1}$在英译可见、法文修改处难辨，照存并注明类型风险。划去说明来自英译编者，既不属于作者数学正文，亦不表示划去部分的证明已补全。
\section*{第二辑：英译物理第68--73页}
\paragraph{第56--57批：划去稿的定义及命题8。}
辨清(a)不是(a)+(b)；核由指定次数的分量所生成不能简化为仅该分量非零。命题8及附带引理在划去稿中有秩计数障碍。即使修正计数，秩下界仍不推出二次关系存在；取完全自由的对称代数即可反驳这种推法。这里只恢复扫描可读部分，不把草稿的缺步改为有证命题。
\paragraph{第58--59批：命题9与推论10。}
分辨英译段次和商式、原稿的自然二次类与提升不变性。推论中恢复法文可以确认的条件，英译漏条件的反例另列，保留划去边界；证明重建是译校说明，不属于作者原证明。
\paragraph{第60--61批：导秩分层与二次初始项。}
法文的$x\notin V''$非英译$x\in V''$，否则落在恰须避开的降秩层；明确相对微分基底$Y$。设余核秩为$r=d-m+1$，行列式张量的可逆因子为$K^{-r}$，原文把指数写作$-n+m-1$。核与余核的局部自由说法须在有恒秩的概形层上核对，单知各剩余域的秩不够。二次初始项的零超曲面光滑也不等同于所有特征下的极化判别式非退化。
\paragraph{第62批：光滑塔、分歧、余维。}
英译误把$d$归给$Y$，与$N=\dim(Y/S)$及$N+d$冲突；正文明确$d=\dim(X/Y)$。法英均将最大子式降秩簇的余维下界写错：$(t,0)$的例子余维$1$而$d-m+1=2$，通常只有分支余维上界。按前$m-1$行的单位子式建立 étale 局部坐标，可把分歧理想写作相对微分模的零阶Fitting理想并由 Hessian 行列式生成。该补证在明确取此源端分歧结构时成立，不是扫描稿中已有的证明，也没有在特征$2$除以$2$。
\section*{第三辑：英译物理第73--77页}
\paragraph{第63批：IV.17.15题名与术语。}
法文物理112页手改题名为“光滑型：增广与初等奇点”，并非全节只论二次型。法英都漏掉“超奇异”叙述的一重否定，紧接一句说明核不能由非退化型生成，故改正文并注原字。特征$2$时$x^2+yz$的射影零曲线光滑而极化有根基，术语的逆推须加偶数秩或特征不为$2$。
\paragraph{第64批：草稿勘误的析取。}
法文物理113页是“且”，英译是“或”；这两层并存于正文脚注。Katz \href{https://grothendiecksga.com/read/sga7/fr/17-3.html}{SGA 7 II，XVII §3} 的 Hessian 判据支持在其光滑射影场景的析取，不是对草稿整条IV.23.9.2的证明。第$n$次嵌入与$f_n=u_n\circ f$的对应属交叉引用推断，特征$2$的不可分次数猜想仍是猜想。
\paragraph{第65批：15.1的一般光滑型。}
正秩纤维的“不等于整个射影空间”才等同于型非零；秩零不应把非零常量当成非空射影空间上的条件。在任意特征使用$\sum X_i\partial_i\varphi=n\varphi$不要求除法；当$n$在剩余域为零，不能把方程$\varphi$从雅可比理想的生成元中无条件删去。环表述的省略仍保留。
\paragraph{第66批：15.2的初等增广及几何定义。}
余法模明定为$J/J^2$；法文的下标顺序优先。初始齐次关系在次数$n$形成局部可逆直和因子，替换生成元的单位倍数不换射影零概形。按极大理想分次商计算可得 Hilbert--Samuel 重数$n$（含零维）。比较不同的有理提升须在剩余域的共同扩域中进行，不从任意域扩张推出绝对与几何无条件相同。
\paragraph{第67批：15.3--15.6的局部刻画。}
SUP为上文，不是英译排出的并号表达式；§15.5只指示如何定义。由法文上下文恢复$X\subseteq X''\subseteq X'$。正则局部环境可先删线性参数，使初始理想由一个型生成，继而以初始项消去及Krull交推出主理想；一般诺特环的完全交用完备化解释。吹起后例外纤维为正则有效Cartier除子，利用恰当性得到穿孔谱正则，再用Serre判据核对约化与正规维数阈值。零维可用Cohen结构得DVR之商，不限等特征。
\paragraph{第68批：15.7的几何变式。}
法文分开$d=\dim_xX$与$n$，右边手加“几何”，左边也据数学范围补明。厚对角的截断主部末层是$I^n/I^{n+1}$，不取$n!$为分母；对规范有理提升先判次$n$非零光滑型。孤立性用非光滑 locus 基变与概形在点处维数不变，不能换成一般非闭点的局部环维数。有限可分情形规范提升延拓极大理想；删掉“有限”的猜测被节点柱面泛点反例否定。纯不可分单点又说明几何初等不推出绝对初等。下一断点是附录IV.17.16光滑二次型，英物理78、法物理161；后文的二次型判据不因此已校译。
\section*{仍需继续的核查}
旧批附件所记“未合并”是历史状态，以第三辑接续记录为准。保留旧卷的残文及未决判断；局部校订不等同于全卷或全部证明的复审。
